\chapter{Entrega e Arquivamento}

A entrega do TCC deve seguir as orientações do curso e da instituição. Além do conteúdo acadêmico, é necessário observar formato do arquivo, prazos, versão correta, identificação e eventuais documentos complementares.

\section{Versão para banca}

A versão encaminhada à banca deve estar completa, revisada e normalizada. Não é adequado enviar texto inacabado, com referências incompletas, seções sem revisão ou elementos obrigatórios ausentes.

Antes de enviar, confira:

\begin{itemize}
  \item capa e folha de rosto;
  \item resumo e palavras-chave;
  \item sumário atualizado;
  \item títulos e numeração;
  \item citações e referências;
  \item figuras, tabelas e fontes;
  \item paginação;
  \item anexos e apêndices citados;
  \item ortografia e clareza.
\end{itemize}

\section{Versão final}

A versão final é a versão corrigida após avaliação, quando houver banca ou parecer. Ela deve incorporar ajustes exigidos e passar por nova conferência formal.

O estudante deve evitar múltiplos arquivos com nomes confusos. Recomenda-se manter controle claro de versões, com data e indicação de finalidade.

\section{Formato do arquivo}

O formato de entrega deve seguir a orientação institucional. Em geral, o PDF é o formato adequado para preservar paginação, fontes, imagens e apresentação final.

Arquivos editáveis podem ser solicitados pelo curso, mas não substituem o PDF da versão destinada à entrega quando a apresentação precisa ser preservada.

\section{Arquivamento}

O estudante deve guardar cópias do trabalho final, fontes utilizadas, instrumentos, dados autorizados, pareceres, certificados e comprovantes de submissão. O arquivamento responsável facilita comprovação, revisão futura e continuidade acadêmica.

\section{Publicação e divulgação}

Quando o TCC for publicado em repositório, evento, artigo ou outro meio, devem ser observados direitos autorais, autorização de participantes, sigilo, licenças de imagem e eventuais restrições institucionais.

Resultados de pesquisa não devem ser divulgados de forma que exponha participantes ou instituições sem autorização.
